% META: {"id": "TNSI-LANG-AM-EXTRAIT", "chapitre": "TNSI-LANGAGES-ET-PROGRAMMATION", "type_objet": "amenagee", "capacites_codes": ["C4", "C5", "C11"], "derive_de": ["TNSI-LANG-EX-001", "TNSI-LANG-EX-004"], "profil_amenagement": ["SEQUENCAGE", "ALLEGEMENT_ENONCE", "REPONSE_FERMEE", "AERATION", "AMORCE_FOURNIE", "RETRAIT_DOUBLE_TACHE"], "mode_creation": "adaptation", "sources_inspiration": [], "status": "generated"}
\section*{Version aménagée --- Extrait Langages et programmation}
\textit{Consignes séquencées, mise en page aérée, énoncés allégés.}

\bigskip

\subsection*{Exercice 1 --- Une fonction récursive (C4, C5)}

On veut la somme des chiffres d'un entier positif. Exemple : $327 \rightarrow 12$.

\bigskip

\textbf{Étape 1.} Décompose $327$. Une ligne à la fois.

\medskip
\textit{Rappel :} \lstinline{n % 10} donne le dernier chiffre,
\lstinline{n // 10} enlève ce dernier chiffre.

\bigskip

\begin{center}
\begin{tabular}{|c|c|c|}
\hline
$n$ & \textbf{Dernier chiffre} & \textbf{Reste} \\
\hline
$327$ & $7$ & $32$ \\
\hline
$32$ & \dotfill & \dotfill \\
\hline
$3$ & \dotfill & \dotfill \\
\hline
\end{tabular}
\end{center}

\bigskip

\textbf{Étape 2.} Quand la décomposition s'arrête-t-elle ?

\medskip
Quand $n$ vaut \dotfill

\bigskip

\textbf{Étape 3.} Complète la fonction. Le cas d'arrêt d'abord.

\begin{python}
def somme_chiffres(n):
    if n == ______:
        return 0
    return n % 10 + somme_chiffres(n ______ 10)
\end{python}

\medskip
Le second trou : \fbox{\texttt{//}} \qquad ou \qquad \fbox{\texttt{\%}}

\bigskip

\textbf{Étape 4.} Sans cas d'arrêt, que se passerait-il ?

\medskip
\begin{center}
\fbox{la fonction s'appelle sans fin} \qquad \fbox{elle renvoie $0$}
\end{center}

\bigskip
\bigskip

\subsection*{Exercice 2 --- Un effet de bord (C11)}

\begin{python}
def ajouter_note(bulletin, note):
    bulletin.append(note)
    return bulletin

mien = [12, 15]
tien = ajouter_note(mien, 18)
\end{python}

\bigskip

\textbf{Étape 1.} Complète le tableau après ces quatre lignes.

\bigskip

\begin{center}
\begin{tabular}{|c|c|}
\hline
\textbf{Variable} & \textbf{Valeur} \\
\hline
\texttt{tien} & \texttt{[12, 15, 18]} \\
\hline
\texttt{mien} & \dotfill \\
\hline
\end{tabular}
\end{center}

\bigskip

\textbf{Étape 2.} La liste de départ a-t-elle été modifiée ?

\medskip
\begin{center}
\fbox{Oui} \qquad \fbox{Non}
\end{center}

\medskip
C'est ce qu'on appelle un \dotfill\ de bord.

\bigskip

\textbf{Étape 3.} Complète la version sans effet de bord.

\begin{python}
def ajouter_note(bulletin, note):
    copie = ______(bulletin)
    copie.append(note)
    return copie
\end{python}

\medskip
Le mot manquant : \fbox{\texttt{list}} \qquad ou \qquad \fbox{\texttt{len}}
